\section{Planteamiento del Problema}
    Dado un conjunto de imágenes con diferentes problemas:
    \begin{enumerate}[label=\alph*.]
        \item Bajo contraste.
        \item Sobreexposición.
        \item Subexposición.
        \item Iluminación no uniforme.
        \item Predominio de determinado canal de color.
    \end{enumerate}
    Se deberá que desarrollar un pequeño sistema de análisis y mejoramiento de imágenes, 
    el cual deberá que:
    \begin{enumerate}
        \item Cargar una imagen.
        \item Identificar sus dimensiones y número de canales.
        \item Mostrar histogramas.
        \item Genera versión en escala de grises.
        \item Modificar brillo.
        \item Modificar contraste.
        \item Aplicar transformación gamma.
        \item Realizar ecualización de histograma.
        \item Comparar las imágenes antes/después.
        \item Generar un reporte de resultados.
    \end{enumerate}

    \subsection{Restricciones}
        Al menos dos transformaciones deberán implementarse manualmente con NumPy, sin utilizar 
        una función de alto nivel que haga directamente todo el procedimiento. Por ejemplo, no 
        basta con llamar a una función de ecualización y entregar el resultado.